% ============================================================
%  Modul 10 — Implementasi Mutex dan Locks
% ============================================================

\chapter{Implementasi Mutex dan Locks}
\label{chap:modul-10}

% --- 1. Tujuan dan Output Praktikum ---
\section{Tujuan dan Output Praktikum}

\subsection{Tujuan Praktikum}
Setelah menyelesaikan praktikum ini, mahasiswa diharapkan mampu:
\begin{enumerate}
	\item Memahami cara kerja \textit{semaphore} sebagai mekanisme dasar sinkronisasi antar \textit{thread}.
	\item Memahami konsep \textit{mutex} (\textit{Mutual Exclusion Lock}) dan perbedaannya dengan \textit{semaphore} secara semantik maupun struktural.
	\item Mengimplementasikan mekanisme penahanan (\textit{wait}) dan pelepasan (\textit{signal}) sumber daya, serta sistem penguncian eksklusif pada tingkat kernel.
	\item Memverifikasi efektivitas mekanisme sinkronisasi yang dibangun dalam mencegah dan menghilangkan \textit{race condition} yang diamati pada modul sebelumnya.
\end{enumerate}

\subsection{Output Praktikum}
Pada akhir praktikum ini, mahasiswa diharapkan menghasilkan:
\begin{itemize}
	\item Implementasi algoritma sinkronisasi kernel yang valid dan mampu meregulasi antrean \textit{thread} dengan tepat.
	\item Simulasi sistem operasi yang mendemonstrasikan penyelesaian skenario \textit{Producer-Consumer} dengan eksekusi yang tersinkronisasi, rapi, dan konsisten.
	\item Laporan analisis yang menjabarkan bagaimana struktur \textit{semaphore} dan \textit{lock} bekerja sama secara harmonis untuk melindungi area \textit{critical section} dari interupsi paksa.
\end{itemize}

% --- 2. Dasar Teori ---
\newpage
\section{Dasar Teori}

\subsection{Konsep Semaphore}

\textit{Semaphore} adalah sebuah struktur data abstrak yang digunakan untuk mengoordinasikan akses beberapa \textit{thread} atau proses terhadap sumber daya (\textit{shared resource}) yang jumlahnya terbatas.
Konsep ini pertama kali diperkenalkan oleh ilmuwan komputer Edsger W. Dijkstra pada tahun 1965.

Secara konseptual, \textit{semaphore} bekerja seperti "keranjang berisi token". Jika sebuah \textit{thread} ingin melakukan pekerjaan, ia harus mengambil satu token dari keranjang. Jika keranjang kosong,
\textit{thread} tersebut harus menunggu (mengantre) sampai ada token yang dikembalikan oleh \textit{thread} lain.

Di dalam sistem operasi, sebuah \textit{semaphore} terdiri dari dua komponen utama:
\begin{enumerate}
	\item Sebuah variabel bilangan bulat (integer) non-negatif yang mewakili "jumlah token" atau kapasitas sumber daya.
	\item Sebuah antrean (\textit{waiting queue}) untuk menampung \textit{thread} yang sedang menunggu giliran.
\end{enumerate}

Semaphore dikendalikan melalui dua operasi atomik (tidak dapat diinterupsi):
\begin{itemize}
	\item \textbf{Operasi \texttt{Wait(S)} (P atau \textit{Proberen})}: Mengurangi nilai variabel semaphore. Jika nilainya menjadi $0$ (atau kurang), maka \textit{thread} yang mengeksekusi operasi ini akan diubah
	      statusnya menjadi \textit{Blocked} / \textit{Asleep} dan dimasukkan ke dalam antrean. Pada sistem OS, operasi ini dinamakan \texttt{sema\_down()}.
	\item \textbf{Operasi \texttt{Signal(S)} (V atau \textit{Verhogen})}: Menambah nilai variabel semaphore. Jika ada \textit{thread} yang sedang mengantre (tertidur), operasi ini akan membangunkan (\textit{Wake up})
	      salah satu \textit{thread} tersebut agar siap dieksekusi kembali oleh CPU. Pada sistem OS, operasi ini dinamakan \texttt{sema\_up()}.
\end{itemize}

Berdasarkan kapasitasnya, semaphore dibagi menjadi dua:
\begin{enumerate}
	\item \textbf{Counting Semaphore}: Nilainya bisa berapa saja ($\geq 0$). Sangat berguna untuk mengontrol akses ke *resource* yang memiliki batasan jumlah, seperti jumlah tempat parkir kosong, atau sisa kapasitas \textit{buffer}.
	\item \textbf{Binary Semaphore}: Nilainya hanya dibatasi pada angka $0$ dan $1$. Sering digunakan untuk menjamin \textit{Mutual Exclusion}.
\end{enumerate}

\subsection{Atomisitas Operasi Sinkronisasi}
Meskipun \textit{semaphore} terlihat seperti bilangan pencacah biasa, operasi \texttt{Wait} dan \texttt{Signal} tidak boleh diperlakukan sebagai aritmetika biasa yang bisa disela sembarang waktu. Pada tingkat kernel, langkah "memeriksa nilai, mengubah nilai, lalu memutuskan apakah \textit{thread} harus tidur atau dibangunkan" harus dianggap sebagai satu kesatuan atomik. Jika tidak, dua \textit{thread} dapat membaca keadaan lama yang sama lalu mengambil keputusan yang saling bertabrakan.

Salah satu kegagalan klasik adalah \textit{lost wakeup}. Misalnya, sebuah \textit{thread} melihat nilai semaphore masih nol dan baru akan tidur, tetapi tepat sebelum benar-benar masuk antrean, \textit{thread} lain sudah lebih dulu melakukan \texttt{Signal}. Jika operasi itu tidak atomik, sinyal kebangkitan dapat "hilang": satu pihak sudah mengirim sinyal, tetapi pihak lain tetap tertidur dan sistem masuk keadaan tidak konsisten. Inilah alasan implementasi \texttt{sema\_down()} dan \texttt{sema\_up()} di kernel biasanya dilindungi dari interleaving yang berbahaya, misalnya dengan menonaktifkan interupsi pada bagian kritis yang sangat pendek.

\subsection{Status \textit{Thread}: Ready, Running, dan Blocked}
Ketika sinkronisasi diterapkan, \textit{thread} tidak hanya bergantian berjalan, tetapi juga dapat keluar sementara dari persaingan CPU. \textit{Thread} yang siap dipilih \textit{scheduler} berada pada status \textit{Ready}. Saat ia sedang memakai CPU, statusnya \textit{Running}. Jika ia memanggil \texttt{Wait} pada semaphore yang belum tersedia, statusnya berubah menjadi \textit{Blocked} dan \textit{thread} tersebut dipindahkan ke antrean penunggu milik objek sinkronisasi.

Perbedaan \textit{Blocked} dengan \textit{busy waiting} sangat penting. Pada \textit{busy waiting}, \textit{thread} terus mengulang pengecekan sambil tetap menghabiskan siklus CPU. Pada pendekatan \textit{blocked}, \textit{thread} justru dikeluarkan dari antrean siap dan baru kembali ke \textit{ready queue} ketika \texttt{Signal} atau \texttt{Unlock} yang relevan terjadi. Model inilah yang lebih efisien dan dipakai dalam praktikum ini melalui pasangan fungsi \texttt{thread\_block()} dan \texttt{thread\_unblock()}.

INS.POST 99:

\noindent \textbf{Invariant dan Urutan Operasi yang Benar:} Pada \textit{bounded buffer}, nilai \texttt{sema\_empty} dan \texttt{sema\_full} harus selalu merepresentasikan jumlah slot kosong dan slot terisi secara konsisten. Karena itu, urutan operasi tidak boleh dibalik sembarangan. \textit{Producer} wajib memastikan slot kosong tersedia lebih dulu, baru memasuki \textit{critical section}. Setelah data benar-benar ditulis, barulah ia memberi sinyal bahwa jumlah data siap baca bertambah. \textit{Consumer} mengikuti pola kebalikannya: pastikan data tersedia, masuk ke area kritis, baca data, lalu umumkan bahwa satu slot kembali kosong. \textit{Mutex} saja tidak cukup menggantikan semaphore counting, karena penguncian hanya melindungi akses serentak, tetapi tidak menghitung kapasitas \textit{buffer}. Jika urutannya salah, sistem dapat menghasilkan jumlah hitung keliru, \textit{deadlock}, atau pembacaan data yang belum valid.

\subsection{Mutex (Mutual Exclusion Lock)}

\textit{Mutex} atau \textit{Lock} adalah mekanisme sinkronisasi yang secara khusus dirancang untuk menjamin \textit{Mutual Exclusion} pada area \textit{Critical Section}.

Analogi paling mudah untuk memahami \textit{Mutex} adalah mengibaratkannya seperti penggunaan bilik kamar mandi.
\begin{itemize}
	\item Ketika seseorang (\textit{Thread A}) memasuki bilik (\textit{Critical Section}), ia akan mengunci pintu dari dalam (\textit{Lock}).
	\item Orang lain (\textit{Thread B}) yang datang kemudian harus mengantre di luar hingga pintu tersebut dibuka (\textit{Unlock}).
	\item Kunci pintu hanya dapat dibuka dari dalam oleh orang yang sama yang sebelumnya menguncinya (\textit{Thread A}).
\end{itemize}

\subsubsection{Konsep Ownership (Kepemilikan)}
Hal ini merupakan perbedaan mendasar antara \textit{Mutex} dan \textit{Semaphore}. \textit{Mutex} memiliki konsep \textbf{Ownership} (Kepemilikan). Sebuah \textit{Mutex} akan mengingat identitas \textit{thread} mana yang berhasil menguncinya. Sistem operasi yang baik
akan menolak operasi \textit{Unlock} jika fungsi tersebut dipanggil oleh \textit{thread} yang bukan pemiliknya.

Di dalam level implementasi (seperti pada OS), Mutex sering kali dibangun di atas \textit{Binary Semaphore} dengan nilai awal $1$, namun ditambah dengan variabel \texttt{holder} untuk mencatat siapa pemegang kunci tersebut.

\begin{table}[h]
	\centering
	\caption{Perbedaan Konseptual Semaphore dan Mutex}
	\begin{tabulary}{\textwidth}{|L|L|L|}
		\hline
		\textbf{Aspek} & \textbf{Semaphore} & \textbf{Mutex (Lock)} \\ \hline
		Analogi & Keranjang berisi token bebas. & Kunci gembok eksklusif. \\ \hline
		Konsep Nilai & Bilangan bulat $\geq 0$ (tergantung kapasitas). & Hanya $0$ (Terkunci) atau $1$ (Bebas). \\ \hline
		Kepemilikan (\textit{Ownership}) & Tidak ada. \textit{Thread A} bisa melakukan \texttt{Wait}, dan \textit{Thread B} bisa melakukan \texttt{Signal}. & Ada. Hanya \textit{thread} pemegang kunci yang berhak melakukan \textit{Release/Unlock}. \\ \hline
		Kegunaan Utama & Sinkronisasi urutan kejadian dan manajemen jumlah sumber daya. & Melindungi \textit{Critical Section} dari modifikasi serentak. \\ \hline
	\end{tabulary}
\end{table}

\subsection{Pendekatan Logika Sinkronisasi Producer-Consumer}

Pada Modul 9, telah didemonstrasikan bahwa masalah \textit{Producer-Consumer} (\textit{Bounded-Buffer}) menyebabkan anomali ketika \textit{buffer} diakses tanpa aturan. Untuk menyelesaikannya, dibutuhkan gabungan dari \textit{Semaphore} dan \textit{Mutex}.

Logika penyelesaiannya membutuhkan 3 instrumen:
\begin{enumerate}
	\item \textbf{\texttt{mutex\_lock}}: Untuk melindungi variabel \textit{buffer} agar tidak ditulis/dibaca bersamaan.
	\item \textbf{\texttt{sema\_empty}}: \textit{Counting Semaphore} untuk menghitung jumlah slot kosong di dalam \textit{buffer}.
	\item \textbf{\texttt{sema\_full}}: \textit{Counting Semaphore} untuk menghitung jumlah slot terisi di dalam \textit{buffer}.
\end{enumerate}

\noindent \textbf{Alur Logika Producer:}
\begin{itemize}
	\item Sebelum memproduksi data, \textit{Producer} harus melakukan \texttt{Wait(sema\_empty)}. Jika kapasitas \textit{buffer} sudah penuh, \textit{thread} ini akan ditidurkan hingga ada slot yang kosong.
	\item Jika ada slot kosong, \textit{Producer} akan mengamankan \textit{buffer} dengan \texttt{Lock(mutex)}.
	\item Setelah data dimasukkan ke dalam \textit{buffer}, \textit{Producer} membuka kembali pengaman dengan \texttt{Unlock(mutex)}.
	\item Terakhir, \textit{Producer} mengirimkan sinyal \texttt{Signal(sema\_full)} untuk memberi tahu \textit{Consumer} bahwa ada data baru yang siap dibaca.
\end{itemize}

\noindent \textbf{Alur Logika Consumer:}
\begin{itemize}
	\item Sebelum membaca data, \textit{Consumer} harus melakukan \texttt{Wait(sema\_full)}. Jika \textit{buffer} sedang kosong, \textit{thread} ini akan ditidurkan hingga ada data baru.
	\item Jika ada data yang tersedia, \textit{Consumer} akan mengamankan \textit{buffer} dengan \texttt{Lock(mutex)}.
	\item Setelah data selesai dibaca, \textit{Consumer} membuka kembali pengaman dengan \texttt{Unlock(mutex)}.
	\item Terakhir, \textit{Consumer} mengirimkan sinyal \texttt{Signal(sema\_empty)} untuk memberi tahu \textit{Producer} bahwa ada slot kosong yang baru saja ditinggalkan.
\end{itemize}

\begin{table}[h]
	\caption{Tabel Ringkasan Dasar Teori Modul 10}
	\label{tab:m10-dasar-teori}
	\centering
	\begin{tabulary}{\textwidth}{|L|L|}
		\hline
		\textbf{Konsep} & \textbf{Penjelasan Singkat} \\ \hline
		Wait / P() / sema\_down
		& Operasi menahan/menurunkan \textit{resource}. Memblokir \textit{thread} jika \textit{resource} habis. \\ \hline
		Signal / V() / sema\_up
		& Operasi melepaskan/menambah \textit{resource}. Membangunkan \textit{thread} dari antrean. \\ \hline
		Ownership
		& Konsep kepemilikan pada \textit{Mutex} yang memastikan hanya sang pengunci yang boleh membuka kunci. \\ \hline
		Thread Blocking
		& Keadaan di mana \textit{thread} diberhentikan paksa oleh OS dan dikeluarkan dari penjadwalan CPU hingga \textit{event} sinkronisasi terjadi. \\ \hline
	\end{tabulary}
\end{table}


% --- 3. Alat dan Bahan ---
\newpage
\section{Alat dan Bahan}
\label{sec:m10-alat-bahan}

\subsection{Perangkat Keras}
\begin{itemize}
	\item PC atau Laptop dengan arsitektur prosesor basis \textit{host} x86/x64 yang mendukung kelanjutan kompilasi emulator.
	\item Memori minimal 4GB RAM ruang eksekusi.
\end{itemize}

\subsection{Perangkat Lunak (Bahan Praktikum)}
\begin{itemize}
	\item OS \textit{Host} utama (Windows dengan WSL/GitBash, Linux, atau macOS).
	\item Text Editor (misalnya Visual Studio Code) untuk menyunting berkas kode bahasa C dan \textit{Assembly}.
	\item Instalasi Docker Engine (\texttt{docker.io}) untuk menyalakan lingkungan isolasi kompilasi.
	\item Emulator Hardware x86: \texttt{qemu-system-x86} (atau instalasi QEMU setara) yang akan mengeksekusi binari hasil OS.
	\item Direktori \textit{Starter Code} OS 32-bit yang sudah disiapkan sebelumnya (berisi direktori \texttt{src} dengan kerangka lengkap \texttt{kernel.c}, \texttt{thread.h}, \texttt{thread.c}, \texttt{switch.asm}, \texttt{scheduler.h}, dsb).
\end{itemize}

% --- 4. Langkah-Langkah Praktikum ---
\newpage
\section{Langkah-Langkah Praktikum}

\subsection{Persiapan Lingkungan}
\label{sec:m10-persiapan}

Pada praktikum ini, Anda akan menuliskan implementasi fungsi sinkronisasi di dalam \texttt{src/synch.c} dan menghubungkannya dengan simulasi di \texttt{src/kernel.c}.

\begin{enumerate}
	\item Buka terminal, pastikan Anda berada di dalam direktori \textit{source code} OS.
	\item Bersihkan hasil \textit{build} sebelumnya agar kompilasi berikutnya menggunakan kode terbaru yang akan Anda tulis.
	      \begin{lstlisting}[language=bash, caption={Bersihkan build sebelumnya}]
make clean
    \end{lstlisting}
\end{enumerate}

% -----------------------------------------------------------------
\subsection{Sesi 1: Mengimplementasikan Logika Semaphore}
\label{sec:m10-sesi-1}

Semaphore adalah fondasi utama dari mekanisme sinkronisasi di OS. Buka file \texttt{src/synch.c} menggunakan teks editor Anda dan cari blok fungsi \texttt{sema\_down} dan \texttt{sema\_up}.

\begin{enumerate}
	\item \textbf{Implementasi \texttt{sema\_down()} (Operasi Wait / P)} \\
	      Fungsi ini bertugas menahan \textit{thread} jika \textit{resource} sedang kosong (nilai 0). Ubah fungsi \texttt{sema\_down()} menjadi seperti berikut:

	      \begin{lstlisting}[language=C, caption={Kode implementasi sema\_down()}]
void sema_down(struct semaphore *sema) {
    /* 1. Selama nilai 0, masukkan thread ke antrean dan tidurkan */
    while (sema->value == 0) {
        sema->waiters[sema->waiter_count++] = thread_current();
        thread_block(); // Serahkan CPU ke thread lain
    }
    /* 2. Jika sudah lolos (nilai > 0), kurangi nilai \textit{resource} */
    sema->value--;
}
    \end{lstlisting}

	\item \textbf{Implementasi \texttt{sema\_up()} (Operasi Signal / V)} \\
	      Fungsi ini bertugas menambah nilai \textit{resource} dan membangunkan \textit{thread} yang sedang antri (jika ada). Ubah fungsi \texttt{sema\_up()} menjadi seperti berikut:

	      \begin{lstlisting}[language=C, caption={Kode implementasi sema\_up()}]
void sema_up(struct semaphore *sema) {
    /* 1. Kembalikan/tambah nilai \textit{resource} */
    sema->value++;
    
    /* 2. Jika ada thread yang menunggu, bangunkan yang pertama (FIFO) */
    if (sema->waiter_count > 0) {
        struct thread *t = sema->waiters[0];
        
        // Geser sisa antrean ke depan (menghapus index 0)
        for (uint32_t i = 0; i < sema->waiter_count - 1; i++) {
            sema->waiters[i] = sema->waiters[i + 1];
        }
        sema->waiter_count--;
        
        // Bangunkan thread agar masuk ke ready queue lagi
        thread_unblock(t);
    }
}
    \end{lstlisting}
\end{enumerate}

% -----------------------------------------------------------------
\subsection{Sesi 2: Mengimplementasikan Mutex Lock}
\label{sec:m10-sesi-2}

Setelah \textit{Semaphore} berfungsi, mekanisme ini akan dibungkus menjadi sebuah \textit{Mutex Lock}. Masih di dalam file \texttt{src/synch.c}, cari bagian \texttt{lock\_acquire} dan \texttt{lock\_release}.

\begin{enumerate}
	\item \textbf{Implementasi \texttt{lock\_acquire()}} \\
	      Gunakan \texttt{sema\_down} untuk mengunci \textit{binary semaphore} milik lock, lalu catat siapa pemiliknya.

	      \begin{lstlisting}[language=C, caption={Kode implementasi lock\_acquire()}]
void lock_acquire(struct lock *lock) {
    sema_down(&lock->semaphore);     // Tunggu & ambil kunci
    lock->holder = thread_current(); // Catat pemegang kunci saat ini
}
    \end{lstlisting}

	\item \textbf{Implementasi \texttt{lock\_release()}} \\
	      Hapus kepemilikan lock, lalu gunakan \texttt{sema\_up} untuk membebaskan kunci agar bisa dipakai \textit{thread} lain.

	      \begin{lstlisting}[language=C, caption={Kode implementasi lock\_release()}]
void lock_release(struct lock *lock) {
    lock->holder = NULL;             // Hapus identitas pemegang kunci
    sema_up(&lock->semaphore);       // Lepas kunci
}
    \end{lstlisting}
	\item \textbf{Simpan} file \texttt{synch.c}. Saat ini \textit{backend} sinkronisasi Anda sudah selesai.
\end{enumerate}

% -----------------------------------------------------------------
\subsection{Sesi 3: Integrasi Sinkronisasi pada Simulasi Kernel}
\label{sec:m10-sesi-3}

Mekanisme sinkronisasi yang telah dibuat tidak akan berfungsi jika tidak diintegrasikan ke dalam alur program utama. Oleh karena itu, bukalah file \texttt{src/kernel.c}. Kita akan melindungi area \textit{critical section} dari skenario Producer-Consumer menggunakan instrumen yang baru saja diselesaikan.

\begin{enumerate}
	\item Cari fungsi \texttt{producer(void *aux)}. Hapus tanda komentar (uncomment) pada blok TODO yang memanggil \texttt{sema} dan \texttt{lock} sehingga kodenya menjadi seperti ini:

	      \begin{lstlisting}[language=C, caption={Melindungi Critical Section pada Producer}]
static void producer(void *aux) {
    (void)aux;
    int item = 0;
    while (1) {
        /* 1. TUNGGU SLOT KOSONG & KUNCI GEMBOK */
        sema_down(&sema_empty);
        lock_acquire(&buf_lock);

        /* --- CRITICAL SECTION --- */
        shared_buffer = item;
        thread_yield(); // Jebakan ini sekarang AMAN karena dilindungi lock
        vga_set_color(VGA_COLOR_BRIGHT_GREEN);
        vga_puts("[Producer] menulis item ");
        vga_putint(item++);
        vga_putchar('\n');
        /* --- END CRITICAL SECTION --- */

        /* 2. BUKA GEMBOK & BERI SINYAL ADA BARANG */
        lock_release(&buf_lock);
        sema_up(&sema_full);

        thread_yield();
    }
}
    \end{lstlisting}

	\item Lakukan hal yang sama pada fungsi \texttt{consumer(void *aux)}.

	      \begin{lstlisting}[language=C, caption={Melindungi Critical Section pada Consumer}]
static void consumer(void *aux) {
    (void)aux;
    while (1) {
        /* 1. TUNGGU BARANG ADA & KUNCI GEMBOK */
        sema_down(&sema_full);
        lock_acquire(&buf_lock);

        /* --- CRITICAL SECTION --- */
        vga_set_color(VGA_COLOR_YELLOW);
        vga_puts("[Consumer] membaca item ");
        thread_yield(); // Jebakan ini sekarang AMAN karena dilindungi lock
        vga_putint(shared_buffer);
        vga_putchar('\n');
        /* --- END CRITICAL SECTION --- */

        /* 2. BUKA GEMBOK & BERI SINYAL SLOT KOSONG */
        lock_release(&buf_lock);
        sema_up(&sema_empty);

        thread_yield();
    }
}
    \end{lstlisting}
	\item \textbf{Simpan} file \texttt{kernel.c}.
\end{enumerate}

% -----------------------------------------------------------------
\subsection{Sesi 4: Kompilasi dan Verifikasi Sistem OS}
\label{sec:m10-sesi-4}

\begin{enumerate}
	\item Lakukan kompilasi sistem operasi Anda dengan \textit{flag} sinkronisasi.
	      \begin{lstlisting}[language=bash, caption={Compile dan jalankan dengan sinkronisasi penuh}]
# Pastikan telah menjalankan make clean sebelumnya
make DEMO=-DSYNCHRONIZATION
make run
    \end{lstlisting}

	\item \textbf{Verifikasi Hasil:} Amati keluaran (\textit{output}) di layar QEMU.
	      \begin{itemize}
		      \item \textbf{Sebelumnya (Modul 9):} Teks bertumpuk dan \textit{error} karena interupsi \texttt{thread\_yield()}.
		      \item \textbf{Sekarang (Modul 10):} Setiap \texttt{"[Producer] menulis item N"} selalu diikuti secara rapi oleh \texttt{"[Consumer] membaca item N"}.
	      \end{itemize}
	      Meskipun CPU dipaksa berpindah \textit{thread} di tengah jalan oleh jebakan \texttt{thread\_yield()}, sistem operasi Anda kini sudah memiliki pengamanan memori.

	\item Ambil tangkapan layar (\textit{screenshot}) hasil eksekusi untuk membuktikan bahwa implementasi sinkronisasi telah berjalan sesuai ekspektasi.
\end{enumerate}

\begin{figure}[h]
	\centering
	\includegraphics[width=0.85\textwidth]{Figure/10/1.jpeg}
	\caption{Keluaran (\textit{output}) QEMU setelah implementasi sinkronisasi penuh --- Producer dan Consumer bergantian secara rapi, Mutex berhasil mengatasi jebakan Preemption.}
	\label{fig:m10-sesi4}
\end{figure}

% ----------------------------------------------------------------
\section{Latihan}

\subsection{Tujuan}

Kerjakan tugas-tugas berikut secara berurutan mengikuti alur sesi pada Modul 9 dan Modul 10. Sertakan bukti berupa tangkapan layar (\textit{screenshot}) QEMU atau penggalan kode yang relevan pada setiap langkah yang diminta.

% ============================================================
\subsection{Sesi 1 — Menjalankan Demo Race Condition}
% ============================================================

\subsubsection{Persiapan dan Kompilasi}

Pastikan implementasi Minggu 4 sudah berfungsi, lalu kompilasi dengan \textit{flag} sinkronisasi diaktifkan.

\begin{enumerate}

	\item Jalankan perintah berikut untuk mengompilasi dan menjalankan kernel dengan \textit{flag} \texttt{SYNCHRONIZATION}:
	      \begin{lstlisting}[language=bash, caption={Compile dan jalankan demo race condition}]
make clean
make DEMO=-DSYNCHRONIZATION
make run
    \end{lstlisting}

	      Amati keluaran (\textit{output}) yang muncul selama $\pm$15 detik, lalu ambil tangkapan layar (\textit{screenshot}) pada layar QEMU. Tempel hasil tangkapan layar di sini:

	      % \begin{figure}[htbp]
	      %     \centering
	      %     % Ganti dengan screenshot output QEMU saat race condition
	      %     \includegraphics[width=0.85\textwidth]{Figure/handson05/1-race.png}
	      %     \caption{Output QEMU Modul 9 --- tanpa sinkronisasi.}
	      %     \label{fig:h05-1-race}
	      % \end{figure}


	\item Perhatikan output yang muncul. Isilah tabel berikut berdasarkan anomali yang kamu temukan pada layar QEMU.

	      \begin{table}[htbp]
		      \centering
		      \caption{Anomali yang ditemukan pada output race condition}
		      \begin{tabulary}{\textwidth}{|c|L|L|}
			      \hline
			      \textbf{No.} & \textbf{Anomali yang Terlihat} & \textbf{Kemungkinan Penyebab} \\ \hline
			      1 & \textit{[contoh: baris terpotong di tengah]} & \textit{[Isi jawaban di sini]} \\ \hline
			      2 & \textit{[Isi jawaban di sini]} & \textit{[Isi jawaban di sini]} \\ \hline
			      3 & \textit{[Isi jawaban di sini]} & \textit{[Isi jawaban di sini]} \\ \hline
		      \end{tabulary}
	      \end{table}

\end{enumerate}

% ────────────────────────────────────────────────────────────
\subsubsection{Analisis Kode Penyebab Race Condition}

\begin{enumerate}

	\item Buka file \texttt{src/kernel.c} dan perhatikan fungsi \texttt{producer()} dan \texttt{consumer()}. Isilah tabel berikut untuk mengidentifikasi letak jebakan \texttt{thread\_yield()} pada masing-masing fungsi.

	      \begin{table}[htbp]
		      \centering
		      \caption{Letak \texttt{thread\_yield()} yang berbahaya pada \textit{critical section}}
		      \begin{tabulary}{\textwidth}{|l|L|L|}
			      \hline
			      \textbf{Fungsi} & \textbf{Posisi thread\_yield()} & \textbf{Akibat jika CPU berpindah di sini} \\ \hline
			      \texttt{producer()} & Setelah \texttt{shared\_buffer = item}, sebelum mencetak log & \textit{[Isi jawaban di sini]} \\ \hline
			      \texttt{consumer()} & Setelah mencetak \texttt{"membaca item "}, sebelum \texttt{putint()} & \textit{[Isi jawaban di sini]} \\ \hline
		      \end{tabulary}
	      \end{table}

	\item Tuliskan satu skenario eksekusi konkret yang menjelaskan mengapa angka yang dicetak \textit{Consumer} bisa berbeda dari yang ditulis \textit{Producer}. Ikuti format langkah demi langkah berikut:

	      \begin{table}[htbp]
		      \centering
		      \caption{Skenario \textit{race condition} langkah demi langkah}
		      \begin{tabulary}{\textwidth}{|c|l|L|}
			      \hline
			      \textbf{Langkah} & \textbf{Thread Aktif} & \textbf{Yang Terjadi} \\ \hline
			      1 & \textit{Producer} & \texttt{shared\_buffer = 3} \\ \hline
			      2 & \textit{Producer} & \texttt{thread\_yield()} dipanggil --- CPU berpindah \\ \hline
			      3 & \textit{Consumer} & \textit{[Isi jawaban di sini]} \\ \hline
			      4 & \textit{Consumer} & \textit{[Isi jawaban di sini]} \\ \hline
			      5 & \textit{Producer} & \textit{[Isi jawaban di sini]} \\ \hline
			      6 & \textit{Consumer} & \textit{[Isi: angka yang tercetak = ?]} \\ \hline
		      \end{tabulary}
	      \end{table}

\end{enumerate}

% ============================================================
\subsection{Sesi 2 — Implementasi \textit{Semaphore}}
% ============================================================

\subsubsection{Struktur \textit{Semaphore} pada implementasi OS}

\begin{enumerate}

	\item Perhatikan definisi \texttt{struct semaphore} di \texttt{src/synch.h}. Isilah tabel berikut berdasarkan kode sumber.

	      \begin{table}[htbp]
		      \centering
		      \caption{Field-field pada \texttt{struct semaphore}}
		      \begin{tabulary}{\textwidth}{|l|l|L|}
			      \hline
			      \textbf{Field} & \textbf{Tipe Data} & \textbf{Fungsi} \\ \hline
			      \texttt{value}         & \texttt{unsigned}           & \textit{[Isi jawaban di sini]} \\ \hline
			      \texttt{waiters[]}     & \texttt{struct thread *[]}  & \textit{[Isi jawaban di sini]} \\ \hline
			      \texttt{waiter\_count} & \texttt{uint32\_t}          & \textit{[Isi jawaban di sini]} \\ \hline
		      \end{tabulary}
	      \end{table}

	\item Perhatikan nilai awal \textit{semaphore} yang diinisialisasi di \texttt{kernel\_main()} pada \texttt{kernel.c}:
	      \begin{lstlisting}[language=C, caption={Inisialisasi \textit{semaphore} di \texttt{kernel\_main()}}]
sema_init(&sema_empty, 1);
sema_init(&sema_full,  0);
    \end{lstlisting}

	      Isilah tabel berikut untuk menjelaskan mengapa nilai awal kedua \textit{semaphore} berbeda.

	      \begin{table}[htbp]
		      \centering
		      \caption{Alasan nilai awal \textit{semaphore}}
		      \begin{tabulary}{\textwidth}{|l|c|L|}
			      \hline
			      \textbf{Semaphore} & \textbf{Nilai Awal} & \textbf{Alasan} \\ \hline
			      \texttt{sema\_empty} & 1 & \textit{[Jelaskan mengapa bernilai 1]} \\ \hline
			      \texttt{sema\_full}  & 0 & \textit{[Jelaskan mengapa bernilai 0]} \\ \hline
		      \end{tabulary}
	      \end{table}

\end{enumerate}

% ────────────────────────────────────────────────────────────
\subsubsection{Implementasi sema\_down() dan sema\_up()}

\begin{enumerate}

	\item Salin kode yang telah diimplementasikan di sini:

	      \begin{lstlisting}[language=C, caption={Implementasi \texttt{sema\_down()} oleh mahasiswa}]
void sema_down(struct semaphore *sema) {

    /* Tulis implementasi Anda di sini */

}
    \end{lstlisting}

	      \begin{lstlisting}[language=C, caption={Implementasi \texttt{sema\_up()} oleh mahasiswa}]
void sema_up(struct semaphore *sema) {

    /* Tulis implementasi Anda di sini */

}
    \end{lstlisting}

	\item Jelaskan mengapa \texttt{sema\_down()} menggunakan \texttt{while} dan bukan \texttt{if} untuk mengecek \texttt{sema->value == 0}. Apa yang dapat terjadi jika menggunakan \texttt{if}?

\end{enumerate}

% ============================================================
\subsection{Sesi 3 — Implementasi \textit{Lock} dan Integrasi Kernel}
% ============================================================

\subsubsection{Implementasi lock\_acquire() dan lock\_release()}

\begin{enumerate}

	\item Masih di \texttt{src/synch.c}, implementasikan \texttt{lock\_acquire()} dan \texttt{lock\_release()}. Salin kode yang telah diimplementasikan di sini:

	      \begin{lstlisting}[language=C, caption={Implementasi \texttt{lock\_acquire()} oleh mahasiswa}]
void lock_acquire(struct lock *lock) {

    /* Tulis implementasi Anda di sini */

}
    \end{lstlisting}

	      \begin{lstlisting}[language=C, caption={Implementasi \texttt{lock\_release()} oleh mahasiswa}]
void lock_release(struct lock *lock) {

    /* Tulis implementasi Anda di sini */

}
    \end{lstlisting}

	\item Isilah tabel berikut untuk menjelaskan peran \textit{field} \texttt{holder} pada \texttt{struct lock}.

	      \begin{table}[htbp]
		      \centering
		      \caption{Peran \textit{field} \texttt{holder} pada \texttt{struct lock}}
		      \begin{tabulary}{\textwidth}{|L|L|}
			      \hline
			      \textbf{Pertanyaan} & \textbf{Jawaban} \\ \hline
			      Apa yang disimpan di \texttt{holder}?
			      & \textit{[Isi jawaban di sini]} \\ \hline
			      Kapan \texttt{holder} diisi (\textit{non-NULL})?
			      & \textit{[Isi jawaban di sini]} \\ \hline
			      Kapan \texttt{holder} dikosongkan (\texttt{NULL})?
			      & \textit{[Isi jawaban di sini]} \\ \hline
			      Apa akibatnya jika \texttt{holder} tidak pernah dibersihkan (\textit{reset}) ke \texttt{NULL}?
			      & \textit{[Isi jawaban di sini]} \\ \hline
		      \end{tabulary}
	      \end{table}

\end{enumerate}

% ────────────────────────────────────────────────────────────
\subsubsection{Aktivasi Sinkronisasi dan Analisis \textit{Critical Section}}

\begin{enumerate}

	\item Buka \texttt{src/kernel.c}. Hilangkan tanda komentar (\textit{uncomment}) pada keempat blok \texttt{TODO} di fungsi \texttt{producer()} dan \texttt{consumer()} untuk mengaktifkan pengaman.
	      Perhatikan bahwa kini area manipulasi data telah diapit oleh fungsi sinkronisasi (membentuk \textit{Entry Section} dan \textit{Exit Section}).

	      \begin{lstlisting}[language=C, caption={Tampilan \texttt{producer()} setelah hilangkan blok komentar}]
static void producer(void *aux) {
    (void)aux;
    int item = 0;
    while (1) {
        sema_down(&sema_empty);   /* <- sudah di-uncomment */
        lock_acquire(&buf_lock);  /* <- sudah di-uncomment */

        /* --- CRITICAL SECTION --- */
        shared_buffer = item;
        thread_yield();
        vga_puts("[Producer] menulis item ");
        vga_putint(item++);
        vga_putchar('\n');
        /* --- END CRITICAL SECTION --- */

        lock_release(&buf_lock);  /* <- sudah di-uncomment */
        sema_up(&sema_full);      /* <- sudah di-uncomment */
        thread_yield();
    }
}
    \end{lstlisting}

	\item Secara teoritis, \textit{Mutual Exclusion} (\textit{Mutex}) digunakan untuk melindungi variabel global di memori (dalam hal ini \texttt{shared\_buffer}). Namun pada kode di atas,
	      instruksi pencetakan ke layar (\texttt{vga\_puts} dan \texttt{vga\_putint}) juga ikut dimasukkan ke dalam area \textit{Critical Section} yang digembok oleh \texttt{buf\_lock}.

	      Analisislah mengapa perangkat antarmuka/layar (VGA) juga dikategorikan sebagai \textit{shared resource} dalam sebuah Sistem Operasi! Apa yang akan terjadi pada integritas data dan
	      tampilan di layar jika \texttt{lock\_release(\&buf\_lock)} dipindahkan ke atas, tepat setelah baris \texttt{shared\_buffer = item;} (sebelum mencetak log)?

\end{enumerate}

% ============================================================
\subsection{Sesi 4 — Kompilasi dan Verifikasi Sistem}
% ============================================================

\begin{enumerate}

	\item \textit{Compile} ulang dan jalankan kernel dengan sinkronisasi penuh:
	      \begin{lstlisting}[language=bash, caption={Compile dan jalankan dengan sinkronisasi penuh}]
make clean
make DEMO=-DSYNCHRONIZATION
make run
    \end{lstlisting}

	      Sertakan tangkapan layar (\textit{screenshot}) keluaran (\textit{output}) QEMU \textbf{sebelum} dan \textbf{sesudah}  sinkronisasi diterapkan secara berdampingan.

	      % \begin{figure}[H]
	      %     \centering
	      %     \begin{minipage}[t]{0.48\textwidth}
	      %         \centering
	      %         % Ganti dengan screenshot Modul 9 (race condition)
	      %         \includegraphics[width=\textwidth]{Figure/handson05/4-sebelum.png}
	      %         \caption*{\textbf{Sebelum}  --- output kacau.}
	      %     \end{minipage}%
	      %     \hfill
	      %     \begin{minipage}[t]{0.48\textwidth}
	      %         \centering
	      %         % Ganti dengan screenshot Modul 10 (sudah sinkron)
	      %         \includegraphics[width=\textwidth]{Figure/handson05/4-sesudah.png}
	      %         \caption*{\textbf{Sesudah}  --- output rapi.}
	      %     \end{minipage}
	      %     \caption{Perbandingan output QEMU sebelum dan sesudah sinkronisasi.}
	      %     \label{fig:h05-4-compare}
	      % \end{figure}


	\item Verifikasi urutan keluaran (\textit{output}) yang benar. Ekspektasi keluaran (\textit{output}) setelah sinkronisasi berhasil adalah sebagai berikut:

	      \begin{lstlisting}[caption={Ekspektasi keluaran (\textit{output}) \textit{Producer-Consumer} yang tersinkronisasi}]
[Producer] menulis item 0
[Consumer] membaca item 0
[Producer] menulis item 1
[Consumer] membaca item 1
[Producer] menulis item 2
[Consumer] membaca item 2
...
    \end{lstlisting}

	\item Apakah keluaran (\textit{output}) pada QEMU sudah konsisten seperti di atas? Jelaskan mengapa \texttt{thread\_yield()} yang ada di dalam \textit{critical section} tidak lagi menyebabkan anomali setelah \texttt{lock} dipasang.

	\item Isilah tabel ringkasan peran ketiga instrumen sinkronisasi berdasarkan hasil praktikum yang telah dilakukan.

	      \begin{table}[htbp]
		      \centering
		      \caption{Peran tiga instrumen sinkronisasi pada \textit{Producer-Consumer}}
		      \begin{tabulary}{\textwidth}{|l|l|L|}
			      \hline
			      \textbf{Instrumen} & \textbf{Jenis} & \textbf{Peran} \\ \hline
			      \texttt{sema\_empty}
			      & \textit{Counting Semaphore}
			      & \textit{[Jelaskan kapan Producer harus menunggu]} \\ \hline
			      \texttt{sema\_full}
			      & \textit{Counting Semaphore}
			      & \textit{[Jelaskan kapan Consumer harus menunggu]} \\ \hline
			      \texttt{buf\_lock}
			      & \textit{Mutex Lock}
			      & \textit{[Jelaskan apa yang dilindungi]} \\ \hline
		      \end{tabulary}
	      \end{table}

\end{enumerate}

% ============================================================
\subsection{Eksplorasi Mandiri}
% ============================================================

\begin{enumerate}

	% ── Eksplorasi 1 ─────────────────────────────────────────
	\item \textbf{Eksperimen tanpa \textit{Lock}.}\\
	      Pada \texttt{kernel.c}, jadikan blok komentar (\textit{comment out}) untuk pemanggilan \texttt{lock\_acquire()} dan \texttt{lock\_release()} pada \texttt{producer()} dan \texttt{consumer()}, tetapi biarkan \texttt{sema\_down()} dan \texttt{sema\_up()} tetap aktif. \textit{Compile} ulang dan jalankan.

	      \begin{lstlisting}[language=bash, caption={Kompilasi untuk eksperimen tanpa \textit{lock}}]
make clean && make DEMO=-DSYNCHRONIZATION && make run
    \end{lstlisting}

	      \begin{enumerate}
		      \item Sertakan tangkapan layar (\textit{screenshot}) hasil QEMU dari eksperimen ini.

		            % \begin{figure}[htbp]
		            %     \centering
		            %     % Ganti dengan screenshot eksperimen tanpa lock
		            %     \includegraphics[width=0.85\textwidth]{Figure/handson05/exp-tanpa-lock.png}
		            %     \caption{Output QEMU tanpa \texttt{buf\_lock} (semaphore tetap aktif).}
		            %     \label{fig:h05-exp-lock}
		            % \end{figure}


		      \item Apakah keluaran (\textit{output}) masih kacau atau sudah rapi? Jelaskan mengapa.
	      \end{enumerate}

	      % ── Eksplorasi 2 ─────────────────────────────────────────
	\item \textbf{Eksperimen tanpa \textit{Semaphore}.}\\
	      Kembalikan \texttt{lock\_acquire()} dan \texttt{lock\_release()} seperti semula, lalu jadikan blok komentar (\textit{comment out}) untuk pemanggilan semua \texttt{sema\_down()} dan \texttt{sema\_up()}. \textit{Compile} ulang dan jalankan.

	      \begin{enumerate}
		      \item Sertakan tangkapan layar (\textit{screenshot}) hasil QEMU dari eksperimen ini.

		            % \begin{figure}[htbp]
		            %     \centering
		            %     % Ganti dengan screenshot eksperimen tanpa semaphore
		            %     \includegraphics[width=0.85\textwidth]{Figure/handson05/exp-tanpa-sema.png}
		            %     \caption{Output QEMU tanpa semaphore (\texttt{buf\_lock} tetap aktif).}
		            %     \label{fig:h05-exp-sema}
		            % \end{figure}


		      \item Apakah keluaran (\textit{output}) masih konsisten? Apa perbedaan yang dapat diamati dibandingkan saat ketiganya aktif? Jelaskan.
	      \end{enumerate}

	      % ── Eksplorasi 3 ──────────────────────────────
	\item \textbf{Eksperimen Urutan Penguncian yang Keliru (\textit{deadlock}).}\\
	      Kembalikan seluruh kode ke kondisi semula sehingga keluaran (\textit{output}) kembali konsisten sebelum melanjutkan eksperimen ini.\\[0.2cm]
	      Pada fungsi \texttt{producer()} di \texttt{kernel.c}, tukar urutan dua baris berikut sehingga \texttt{lock\_acquire()} dipanggil sebelum \texttt{sema\_down()}:

	      \begin{lstlisting}[language=C, caption={Perbandingan urutan yang benar dan yang diubah}]
/* Urutan yang benar: */
sema_down(&sema_empty);   // (1) periksa ketersediaan slot
lock_acquire(&buf_lock);  // (2) kunci buffer

/* Urutan yang diubah untuk eksperimen ini: */
lock_acquire(&buf_lock);  // (1) kunci buffer terlebih dahulu
sema_down(&sema_empty);   // (2) periksa ketersediaan slot
    \end{lstlisting}

	      Kompilasi ulang dan jalankan.

	      \begin{lstlisting}[language=bash, caption={Kompilasi untuk eksperimen \textit{deadlock}}]
make clean && make DEMO=-DSYNCHRONIZATION && make run
    \end{lstlisting}

	      \begin{enumerate}
		      \item Tempelkan tangkapan layar (\textit{screenshot}) layar QEMU setelah perubahan diterapkan. Catat apakah sistem menghasilkan keluaran (\textit{output}) atau langsung berhenti.

		            % \begin{figure}[htbp]
		            %     \centering
		            %     \includegraphics[width=0.85\textwidth]{Figure/handson05/exp-deadlock.png}
		            %     \caption{Kondisi QEMU akibat urutan penguncian yang keliru.}
		            %     \label{fig:h05-exp-deadlock}
		            % \end{figure}


		      \item Lengkapi tabel alur eksekusi berikut untuk menjelaskan mengapa perubahan urutan tersebut menyebabkan sistem berhenti total.

		            \begin{table}[htbp]
			            \centering
			            \caption{Alur eksekusi akibat urutan penguncian yang keliru}
			            \begin{tabulary}{\textwidth}{|c|l|L|}
				            \hline
				            \textbf{Langkah} & \textbf{Thread Aktif} & \textbf{Yang Terjadi} \\ \hline
				            1 & \textit{Producer}
				            & \texttt{lock\_acquire(\&buf\_lock)} berhasil; \textit{Producer} memegang kunci. \\ \hline
				            2 & \textit{Producer}
				            & \texttt{sema\_down(\&sema\_empty)} dipanggil; nilai \texttt{sema\_empty} = 0, \textit{Producer} diblokir melalui \texttt{thread\_block()}. \\ \hline
				            3 & \textit{Consumer}
				            & Mendapat giliran CPU; \texttt{sema\_down(\&sema\_full)} berhasil. \\ \hline
				            4 & \textit{Consumer}
				            & \texttt{lock\_acquire(\&buf\_lock)} dipanggil --- \textit{[Isi jawaban di sini]} \\ \hline
				            5 & ---
				            & \textit{[Jelaskan mengapa tidak ada thread yang dapat melanjutkan]} \\ \hline
			            \end{tabulary}
		            \end{table}


		      \item Apa kesimpulan yang dapat diambil terkait pentingnya urutan pemanggilan \texttt{sema\_down()} dan \texttt{lock\_acquire()} dalam implementasi sinkronisasi?
	      \end{enumerate}

\end{enumerate}
